\section{Q\&A}\label{QA}
此处回答一些可能会有的问题, 同时也是作为前面内容的补充.

\subsection{我是搞材料的, 学这么多计算机的软件的内容做什么?}
本手册面向的读者群体为希望加入MInDes程序开发, 但软件开发相关领域了解不深的科研工作者. 
为尽可能减轻读者负担, 笔者尽己所能将内容压缩至精简且易于理解的程度. 这些内容或问题在软件开发的过程中会
频繁遇到, 因此具有一定的了解会让开发工作更加顺利.

笔者希望读者能在该手册中快速检索到自己需要的信息, 不奢求逐行看完.

\subsection{你写的一大堆全都是泛泛而谈, 有没有更专业的文档?}
这里提供一份清单, 内容是笔者认为较好的软件官网, 教程指南, 以及别的优秀资源. 前述的链接也被汇总至此.

首先是软件官网下载与安装教程:
\begin{itemize}
    \item \bhref{https://visualstudio.microsoft.com/zh-hans/}{Microsoft VS 官网}: 可供下载Microsoft Visual Studio.
    \item \bhref{https://code.visualstudio.com/}{Microsoft VSCode 官网}: 可供下载Microsoft Visual Studio Code.
    \item \bhref{https://learn.microsoft.com/zh-cn/windows/wsl/install}{WSL 官方安装指南}: 可按照微软给出的安装教程安装WSL.
    \item \bhref{https://learn.microsoft.com/zh-cn/windows/wsl/setup/environment}{WSL 官方开发环境设置指南}: 可按照微软给出的安装教程设置WSL开发环境.
    \item \bhref{https://www.paraview.org/download/}{Paraview 官网}: 可供下载Paraview.
    \item \bhref{https://git-scm.com/download/win}{Git 下载页面}: 可下载供Windows平台下运行的Git客户端.
    \item \bhref{https://learn.microsoft.com/zh-cn/windows/terminal/install}{Windows Terminal 官方下载指南}: 可按照微软给出的教程安装Windows Terminal, 实现在一个窗口打开多个Shell.
    \item \bhref{https://learn.microsoft.com/zh-cn/powershell/scripting/install/installing-powershell-on-windows?view=powershell-7.4}{Powershell7 官方下载指南}: 可按照微软给出的教程安装Powershell7.
    \item \bhref{https://github.com/A-moment096/MInDes-LSP}{MInDes-LSP}: 笔者的VSCode插件: MInDes-LSP的GitHub仓库, 可以为MInDes的输入文件提供语法高亮等功能.
\end{itemize}
以下是一些文档和优秀的操作教程, 有官方出品的, 也有个人写作的:
\begin{itemize}
    \item \bhref{https://learn.microsoft.com/zh-cn/visualstudio/windows/?view=vs-2022&preserve-view=true}{VS 文档集}: 可查阅Visual Studio的方方面面.
    \item \bhref{https://code.visualstudio.com/docs}{VSCode 文档集}: 可查阅Visual Studio Code的各项使用方法.
    \item \bhref{https://git-scm.com/docs}{Git 文档集}: Git 的官方文档.
    \item \bhref{https://git-scm.com/book/en/v2}{Pro Git 电子书}: Git官方出品的Git教程电子书, 涵盖了Git的方方面面.
    \item \bhref{https://liaoxuefeng.com/books/git/introduction/index.html}{廖雪峰的Git教程}: 简单鲜活的Git教程, 主要以命令行模式使用Git.
    \item \bhref{https://en.cppreference.com/w/}{cppreference}: 非常好用的C++ API 参考, 记录有大量的C++在不同语言标准下的语言API, 可信度很高.
    \item \bhref{https://cplusplus.com/}{cplusplus}: 同上, 优秀的C++ API 参考, 同时拥有一些教程和文章以供学习.
    \item \bhref{https://wiki.archlinux.org/title/Main_page}{ArchWiki}: Arch Linux团队维护的百科, 除了介绍了Arch Linux外也介绍了许多Linux系统甚至内核的相关知识, 可在需要时查阅.
    \item \bhref{https://linux.die.net/man/}{Linux man pages}: Linux系统中用到的命令, 函数等的用户手册线上全集. 除了该线上版本外, Linux系统中通过命令\cverb|man|也可以调出该文档的离线版本.
    \item \bhref{http://gnu.ist.utl.pt/manual/manual.html}{GNU Manuals Online}: GNU团队出品的开源自由文档, 记录了GNU Project中软件的用户手册, 如bash, GCC等.
    \item \bhref{https://isocpp.org/}{ISO C++ 官网}: C++标准官网. 内含许多有助于开发的信息.
    \item \bhref{https://www.openmp.org/resources/}{OpenMP 学习资源}: OpenMP官网提供的OpenMP相关教程. 
\end{itemize}
以下是一些优秀的工具类网站, 以及优秀的论坛, 在开发过程中可能可以用到:
\begin{itemize}
    \item \bhref{https://quick-bench.com/}{Quick Benchmark}: 一个可以对代码运行速度进行检测的网站, 可以用以对比两段代码的运行速度差距以便优化算法.
    \item \bhref{https://godbolt.org/}{Compiler Explorer}: 一个强大的C++编译器整合网站, 可以从网站中查到不同编译器之间对代码的处理区别, 提供了代码的汇编结果.
    \item \bhref{https://regexr.com/}{RegExr}: 一个可以测试正则表达式结果的网站, 对Linux的部分使用场景或是需要用到正则表达式的场景会有帮助.
    \item \bhref{https://en.wikipedia.org/wiki/Main_Page}{Wikipedia}: 维基百科, 可以简单轻松获得某个词条的基础信息.
    \item \bhref{https://chatgpt.com/}{ChatGPT}: ChatGPT, 强大的AI搜索问答工具.
    \item \bhref{https://stackoverflow.com/}{Stack Overflow}: 著名程序员问答论坛, 近乎任何程序问题都可以从此处寻得专业且切实有效的答案.
    \item \bhref{https://stackexchange.com/}{Stack Exchange}: 上述论坛的母站. 热门的论坛除了Stack Overflow外, 还有Mathematics, Super User, Ask Ubuntu等等, 帖子质量都很高.
    \item \bhref{https://github.com/}{GitHub}: 知名代码托管平台, 拥有海量开源项目.
\end{itemize}
此外, 是一些个人认为有趣的资源或内容:
\begin{itemize}
    \item \bhref{https://blog.brachiosoft.com/posts/git/}{Git的故事}: 讲述了Git的前世今生, 以及GitHub的诞生.
    \item \bhref{https://linux.cn/}{Linux中国}: 一个已经停运的Linux开源社区, 有很多很好的技术贴子, 尤其是Linux相关的帖子.
    \item \bhref{https://www.ruanyifeng.com/blog/}{阮一峰的网络日志}: 著名科技圈周刊, 每周五更新, 支持RSS订阅.
    \item \bhref{https://learn.microsoft.com/en-us/windows/powertoys/}{PowerToys官网}: Microsoft官方出品的Windows平台强大的工具合集.
    \item \bhref{https://www.voidtools.com/}{Everything官网}: 一款强大的Windows端文件搜索工具, 与Windows自带的搜索功能相比搜索内容更快更全.
    \item \bhref{https://www.runoob.com/}{菜鸟教程}: 一个编程学习教程网站. 与上面的教程相比略显浅薄, 但也足以入门.
    \item \bhref{https://liaoxuefeng.com/}{廖雪峰的主页}: 廖雪峰的个人主页, 拥有许多优秀的教程, 如Python教程, Java教程等. 上述的Git教程是其主页内容的一小部分.
\end{itemize}

\subsection{你写的这些长到看不下去, 有没有省流版?}
如果您具有优秀的互联网搜索能力, 且对章节\ref{Concepts}基础概念部分的大多数概念都具有或浅或深的了解, 使用本手册目的仅为快速同步设置好开发环境,
下面是一份您需要注意的内容:
\begin{itemize}
    \item 在Windows上下载并安装VS, VSCode, GitHub Desktop (可选地, 安装\\Powershell7和Windows Terminal, 安装Git替代GitHub Desktop). 安装\\VS时注意选择安装
    {\sffamily{C++桌面开发(Desktop Development with C++)}}的套件.
    \item 使用GitHub Desktop或Git, 将MInDes或MInDes\_DEV的GitHub仓库下载到本地.
    \item 创建一个VS项目, 将下载到本地的MInDes仓库内的文件添加至这个项目中. 请注意, 只添加程序的源文件. 随后为该项目设置C++标准到C++17, 并打开OpenMP的支持.
    \item 将本地仓库中的如下三个文件:\emph{libfftw3-3.dll, libfftw3f-3.dll, libfftw3l-3.dll}复制到项目文件夹下以便链接器寻找到动态链接库
    (可选地, 自行设置动态链接库的寻址). 至此, Windows端的VS配置就完成了.
    \item 根据官方教程安装WSL, 然后在VSCode中安装对应插件以远程连接至WSL. 在VSCode中安装C++与CMake的插件 (可选地, 安装MInDes-LSP).
    \item 在WSL中安装GCC, GDB, CMake, \emph{readline}, 用VSCode打开MInDes或\\MInDes\_DEV的本地仓库, 用CMake尝试运行程序.
    (可选地, 安装fftw3开发用库到本地并修改\emph{CMakeLists.txt}内容以改变fftw库的寻址).
\end{itemize}
完成上述步骤, 您应该可以在一台新组装的电脑上搭建开发MInDes的必要环境. 请注意, 上述步骤中每一步均可能遇到程度不等的问题, 请善用搜索引擎. 遇到的任何问题也欢迎
提出.

\subsection{感觉看完还是云里雾里的\dots}
不得不承认的一点是, 作为实操手册而言, 编写本手册最佳的方式应为一边操作一边记录具体的操作, 以力求读者对照本手册操作时具有足够的可行性.
但笔者的环境已经配好, 推倒重来并不容易, 因此希望读者在使用本手册时能够对照每一步操作, 如有不妥当的步骤请联系我们.

\subsection{这么多页, 一幅图都没有?}
由于本文采用{\LaTeX}写成, 其图片的处理十分麻烦, 而笔者能力着实有限, 因此您看到的早期版本可能是没有或仅含有少量图片的. 在更后期的版本中可能
会加入更多的图片说明, 也可能考虑采用MarkDown重写本手册以便于更新或加入图片.

\subsection{我觉得你写得不好, 我要自己写/修改/增补内容.}
欢迎! 您可以联系笔者以免费获得一份最新版的本手册的{\LaTeX}代码. 该文档如开发到足够成熟, 可能会考虑将之开源, 
届时会放在某个GitHub仓库中. 本手册将持续更新, 还请各位使用者不吝赐教.

\subsection{我看了这个手册, 也搞定了里面讲的配置, 但是还是不知道怎么开始做MInDes开发, 这个不是开发入门手册吗? }
这个手册的初衷其实是为了帮助读者搭建好MInDes开发的环境配置, 但因涉及内容越来越多, 才更名为开发入门手册. 关于实际开发过程的讲解可能会安排培训, 更新本手册
或者出新的, 专注于开发指导的手册.

\endinput